\documentclass[11pt]{article}
\usepackage[margin=1in]{geometry}
\usepackage{amsmath,amssymb,amsthm,booktabs,needspace}
\usepackage[hidelinks]{hyperref}
\setlength{\emergencystretch}{2em}
\newtheorem{proposition}{Proposition}
\newtheorem{theorem}[proposition]{Theorem}
\newcommand{\mr}{\operatorname{mr}}
\newcommand{\rank}{\operatorname{rank}}
\newcommand{\F}{\mathbb F}
\title{Ternary affine tests and new double-cover quotients}
\author{Internal research working proof}
\date{17 November 2026}
\begin{document}
\maketitle

The complete rank-one $\mathrm{GL}_3(\F_3)$ base search on
the original sixteen targets gives no sharp affine certificate.
Every admissible affine lift of every such base is principal.
The search nevertheless produces sixteen nonconstant
representation orbits on eight knots, with full
$\mathrm{GL}_3(\F_3)$ image. They induce
$\mathrm{SL}_3(\F_3)$ quotients of the double-cover groups.
The natural and dual cover modules are now tested directly,
without meridian extension conditions: all have $\dim Z^1=4$,
below the required seven. No additional target closes.

The positive controls are the already-known connected sum of
three trefoils. Explicit order-three and order-two affine
actions detect its meridional rank four. A separate
thirteen-point projective control has the homology gain needed
for the next cover test; target covers of that degree remain
uncomputed. Uniform arguments below are written working
proofs, unformalized and externally unreviewed; worldwide
novelty is not claimed.

\section{The affine and cocycle bounds}

For a representation $\rho:\Gamma\to\mathrm{GL}(V)$ of a knot
group and a meridian $m$ with $\rank(\rho(m)-I)=1$, let
$Z^1_{\rm adm}(\Gamma,V)$ consist of cocycles satisfying
$c(m)\in\operatorname{im}(\rho(m)-I)$. The condition holds on
every conjugate meridian if it holds on one: after
homogenization, the meridian matrix
\[
 \widehat\rho(m)=\begin{pmatrix}\rho(m)&c(m)\\0&1\end{pmatrix}
\]
has rank-one difference from the identity, and this is
conjugacy invariant. In an arc presentation the condition
therefore gives one scalar per meridian image line.

\begin{proposition}
$\mr(K)\ge\dim Z^1_{\rm adm}(\Gamma,V)$. Also, if
$V^\Gamma=0$ and an admissible cocycle is nonprincipal, then
$\widehat\rho$ has zero common fixed space and
$\mr(K)\ge\dim V+1$.
\end{proposition}
\begin{proof}
On any $d$ meridional generators an admissible cocycle is
determined by $d$ values, each in a one-dimensional image.
The evaluation map is injective, so its domain has dimension
at most $d$. For the second assertion, a common fixed vector
$(v,0)$ has $v=0$. A fixed vector $(v,t)$ with $t\ne0$ would
make $c(g)=t^{-1}(v-\rho(g)v)$ principal. Thus the common fixed
space is zero. The fixed hyperplanes of $d$ generating
meridians have intersection of codimension at most $d$,
forcing $d\ge\dim V+1$.
\end{proof}

This is the retained admissible-cocycle mechanism, without
a common bilinear form. Here the affine space has $27$ points;
this is not an assertion that the original binary
$2^k$-point common-form hypothesis is satisfied.

\section{Two explicit known knot controls}

Let $J$ be the connected sum of three positive trefoils. Its
meridian presentation is
\[
 \langle a,b,c,d\mid aba=bab,\ bcb=cbc,\ cdc=dcd\rangle.
\]
Over $\F_3$ take directions
$v_1=e_1,v_2=e_2,v_3=e_3,v_4=(1,1,1)$, and put
$M_i=I_4+(v_i,0)F_i$. The two functional matrices are
\[
 F^{(3)}=\begin{pmatrix}
 0&1&0&0\\2&0&1&0\\2&2&0&0\\1&0&2&1
 \end{pmatrix},\qquad
 F^{(2)}=\begin{pmatrix}
 1&1&0&0\\1&1&1&0\\0&1&1&0\\2&0&2&1
 \end{pmatrix}.
\]
For the first matrix, $F_i(v_i,0)=0$, every $M_i$ has order
three, and adjacent cross-pairings multiply to $-1$.
Expanding the products shows $M_iM_{i+1}M_i=M_{i+1}M_iM_{i+1}$.
For the second, $F_i(v_i,0)=1$, every $M_i$ has order two,
and adjacent cross-pairings multiply to $1$; the same braid
relations hold. Both functional matrices have rank four.
Consequently the common fixed space of their four meridians
is zero. The preceding fixed-hyperplane argument gives
$\mr(J)\ge4$.

In each case the first three translations vanish. For order
three their linear parts generate all $\mathrm{SL}_3(\F_3)$,
of order $5616$; this is checked independently on all $27$
vectors. The fourth linear part belongs to that group and
has nonzero translation. The linear group is transitive on
nonzero vectors, so the affine image contains all translations
and has order $27\cdot5616=151632$. For order two the first
three linear parts generate a group of order $24$. The fourth
linear part lies in it, and its translation orbit spans
$\F_3^3$. Thus this affine image has order $27\cdot24=648$.
Both actions are transitive on $27$ points.

For a literal diagram control, rebase the last two trefoil
factors to $a$ using the inverse conjugators of $ab$ and
$(bc)(ab)$. The braid relations ensure that they take $b,c$
to $a$. The resulting partners retain all four original
generators: the conjugating words already lie in earlier
factors. Glue three standard trefoil PDs along these common
meridian arcs, then renumber the oriented knot walk. The
retained nine-crossing PD has a four-seed propagation trace.
The Wirtinger-number theorem \cite[Theorem 1.3]{wirt} gives
$b(J)\le4$, proving $\mr(J)=b(J)=4$.

Both diagram root spaces have dimension four, and independent
pointwise checks recover every crossing and all three braid
relations. These are known composite knots, not new cases.
All meridian image directions lie in one three-space despite
the zero common fixed space in dimension four; they are
outside an independent-four-direction search.

\section{Exact scope of the ternary base search}

Every nonidentity rank-one difference $A-I=vf$ in
$\mathrm{GL}_3(\F_3)$ has $f(v)=0$ or $1$. The remaining
possibility $2$ makes $A$ singular. Normalize the first
nonzero coordinate of $v$ to $1$. There are $13$ possible
directions. For $f(v)=0$ there are eight nonzero functionals,
giving $104$ transvections of order three. For $f(v)=1$
there are nine functionals, giving $117$ reflections of
order two. Each type is one $\mathrm{GL}_3(\F_3)$ conjugacy
class, by extending a basis adapted to $v$ and $\ker f$.

The group has order $(27-1)(27-3)(27-9)=11232$. The
centralizers of $I+E_{12}$ and $\operatorname{diag}(2,1,1)$
have orders $108$ and $96$. Normalize the first of the four
seed images to the chosen representative and enumerate all
remaining tuples up to its centralizer. The frozen counts are
\begin{center}
\begin{tabular}{rrrr}
\toprule
Meridian order & Class size & Raw normalized tails & Tuple orbits\\
\midrule
3 & 104 & 1,124,864 & 21,632\\
2 & 117 & 1,601,613 & 34,084\\
\bottomrule
\end{tabular}
\end{center}
Both scalar central matrices act trivially by conjugation;
the stabilizers and orbit weights retain that kernel.

Every representative is propagated through the signed
Wirtinger relations of all original sixteen PDs. No proper
image is discarded. With order-three meridians every knot
has only the constant cyclic base, with admissible root
dimension one. With reflections, the knots
\[
 13a_{129},13a_{135},13a_{136},13a_{137},
 13a_{194},13a_{196},13a_{197},13a_{198}
\]
each have the constant base and two nonconstant conjugacy
orbits. Their nonconstant images have order $11232$, hence
are the full general linear group. Their root dimensions
are three. The other eight targets have only the constant
reflection base of root dimension one.

\begin{theorem}
For every rank-one $\mathrm{GL}_3(\F_3)$ meridian base on these
sixteen knots, every admissible affine cocycle is principal.
In particular the homogeneous four-dimensional mechanism has
no zero-common-fixed-space certificate from these bases.
\end{theorem}
\begin{proof}
For each base the independent verifier constructs the map
$w\mapsto(w-A_iw)_i$, expressing its values in the chosen
meridian image lines. Its rank is one for a constant base
and three for each full-image base. The entire admissible
kernel has the same dimension, so it equals the principal
subspace. Every affine lift therefore fixes a point, and
its homogenization fixes the corresponding nonzero vector.
\end{proof}

If the linear meridian image is the identity instead, normal
generation makes the entire linear image trivial. The affine
image is then cyclic translations, since the knot group has
cyclic abelianization. It cannot be transitive on $27$ points
or give a zero common fixed space in homogeneous dimension
four. This accounts for the omitted rank-zero base case.
Bases over extension fields or higher dimensions are not
covered by the finite enumeration.

\section{Reuse through unrestricted double-cover cocycles}

For a reflection base $M_i$, put $g_i=M_iM_0$. In the
orbifold quotient of the knot group with meridian squares
killed, the double-cover kernel has presentation
\[
 G=\langle x_0,\ldots,x_{n-1}\mid x_0=1,
    \ x_ox_a^{-1}x_ox_b^{-1}=1\text{ at each crossing}\rangle.
\]
This is the retained core-group presentation
\cite{core}. It also follows directly by setting
$m_i=x_it$, $t=m_0$, $t^2=1$; the square relations give
$tx_it=x_i^{-1}$, and each involutory crossing gives the
displayed word. Its deck conjugate is a cyclic conjugate
of its inverse, so no second independent family is omitted.

For every full-image base the induced cover image is all
$\mathrm{SL}_3(\F_3)$, since determinant on the knot image
is precisely meridian parity. For both the natural module
$V=\F_3^3$ and its contragredient dual, compute ordinary
$Z^1(G,V)$ using three unrestricted coordinates per
generator. The negative-letter contribution is
$-\rho(x_i)^{-1}c(x_i)$, with no meridian/root constraint.

\Needspace{8\baselineskip}
\begin{proposition}
If $\dim Z^1(G,V)\ge7$ for a three-dimensional module, then
$\rank G\ge3$. With the retained four-bridge upper and
double-cover bound, this proves $\mr(K)=b(K)=4$.
\end{proposition}
\begin{proof}
If $G$ had two generators, a cocycle would be determined by
two vectors and $\dim Z^1(G,V)\le6$. The retained branched-cover
inequality is $\mr(K)\ge1+\rank G$.
\end{proof}

All $32$ natural/dual tests on the sixteen representation
orbits give $V^G=0$, $\dim Z^1=4$ and $\dim H^1=1$.
These modules do not pass the rank gate. This does not prove
that the groups have rank two, or exclude other
$\mathrm{SL}_3(\F_3)$ modules or coefficient fields.
The representation orbits are not claimed to have distinct
quotient kernels.

The positive ordinary control uses the same triple-trefoil
PD and constant reflection base. Its double cover has trivial
action on $\F_3^3$, with $\dim Z^1=9$, detecting rank three.
The core presentation and every ordinary coefficient are
independently reconstructed by evaluating affine words on
the $27$ vectors.

\section{A precise next gate: thirteen-point cover homology}

The new cover quotients act on the thirteen projective lines
of $\F_3^3$. For the inverse image $H$ of a point stabilizer
and any prime $p\ne13$, Shapiro and the split augmentation
module of dimension twelve give
\[
 b_1(H;\F_p)-b_1(G;\F_p)\ge13
          \quad\Longrightarrow\quad\rank G\ge3.
\]
Indeed the augmentation module has no invariant vector, and
two generators give first cohomology dimension at most
$24-12=12$. The same applies to the dual thirteen-line action.
The natural three-dimensional tests above do not address
this twelve-dimensional module.

There is an exact known control. The triple-trefoil cover has
$G_0=C_3*C_3*C_3$. Send its three standard generators to the
first three linear parts of the order-three control. These
generate all $\mathrm{SL}_3(\F_3)$. Each has four fixed
projective points and three cycles of length three, and the
joint action is transitive. In the induced graph of groups
there are $13+3\cdot7=34$ vertices and $3\cdot13=39$ edges.
It has cycle rank six; the twelve fixed-point vertices retain
their $C_3$ stabilizers and all other vertex groups are
trivial. The standard subgroup graph/Kurosh argument
\cite[Chapter I]{serre} gives
\[
 H_0\cong F_6*\mathop{*}_{i=1}^{12}C_3,\qquad
 b_1(H_0;\F_3)-b_1(G_0;\F_3)=18-3=15\ge13.
\]
The finite projective cycle and connectivity inputs are
independently checked. The subgroup structure uses the stated
classical graph-of-groups import. No target thirteen-point
cover homology has yet been computed.

\clearpage
\section{Verification and remaining scope}

The independent base checker filters all $3^9$ matrices,
constructs the full centralizers as permutations of $27$
vectors, and marks every actual orbit in each raw tail cube.
It checks that these orbits are disjoint and omit no tuple,
independently of the producer's successive-stabilizer method.
Geometric paths between undercrossings and an oriented knot
walk reconstruct the signed diagram. The source parser is
used only for checked legacy arc names. A single pass through
the upper trace and pointwise relator tests replace the
producer's iterative early-rejection propagation.

Every root and ordinary coefficient matrix is rebuilt by
giving one affine translation coordinate a time and evaluating
the full word rightmost first. Exact row ranks prove kernel
completeness, and every supplied basis vector is tested.
Both controls, normalized end-to-end searches, image orders
and principal-subspace ranks pass. Ten altered control,
overlapping-orbit, kernel and ordinary-coefficient inputs
reject across the four verification receipts.

The new facts are these exact finite exclusions and explicit
quotients, not a new named knot equality. The original sixteen
targets, other coefficient modules and fields, target cover
ranks, worldwide novelty and full classical MRC remain open.

\begin{thebibliography}{9}
\small
\bibitem{wirt} R. Blair, A. Kjuchukova, R. Velazquez and P.
Villanueva, Wirtinger systems of generators of knot groups,
Commun. Anal. Geom. 28 (2020), Theorem 1.3.
\url{https://archive.intlpress.com/site/pub/files/_fulltext/journals/cag/2020/0028/0002/CAG-2020-0028-0002-a002.pdf}.
\bibitem{core} D. Silver, L. Traldi and S. Williams, Core
groups, J. Knot Theory Ramifications 33 (2024), 2450013.
\url{https://arxiv.org/abs/2311.02048v2}.
\bibitem{serre} J.-P. Serre, Trees, Springer (1980), Chapter I.
\url{https://doi.org/10.1007/978-3-642-61856-7}.
Classical subgroup graph and normal-form theorem imported;
the full book is not independently re-certified here.
\end{thebibliography}
\end{document}
